% ============================================================
% Capitolo 15 — Monorepo, Company Extension e Codebase Legacy
% ============================================================

\chapter{Monorepo, Company Extension e Codebase Legacy}
\label{ch:monorepo-legacy}

I capitoli precedenti hanno seguito TaskFlow API come progetto singolo. Questo capitolo copre tre scenari avanzati che emergono con la complessità reale del lavoro di squadra.

\section{Monorepo: un framework, più progetti}
\label{sec:15-monorepo}

In un monorepo, ogni sottoprogetto che ha bisogno di un contesto indipendente crea il proprio \texttt{\_CONTEXT.md}:

\begin{lstlisting}[language=,caption={Struttura di un monorepo AI-DLC}]
taskflow/
+-- apps/
|   +-- api/           ← REST API
|   |   +-- _CONTEXT.md (Phase 3)
|   +-- worker/        ← Job worker
|   |   +-- _CONTEXT.md (Phase 2)
|   +-- dashboard/     ← Frontend React
|       +-- _CONTEXT.md (Phase 3)
+-- AGENTS.md          ← Framework condiviso
+-- CLAUDE.md
+-- .ai-dlc/             ← Framework condiviso
\end{lstlisting}

L'agente usa il \texttt{\_CONTEXT.md} più vicino al file su cui lavora.

\textbf{Scaffold di un nuovo sottoprogetto:}

\begin{lstlisting}[language=bash,caption={Inizializzazione sottoprogetto in monorepo}]
bash .ai-dlc/tools/init.sh --project-root apps/worker --framework-root .
\end{lstlisting}

\textbf{Aggiornamento indice:} dopo aggiunte o spostamenti di sottoprogetti, aggiorna \texttt{.ai-dlc/projects.json}:

\begin{lstlisting}[language=bash,caption={Aggiornamento indice progetti}]
bash .ai-dlc/tools/update-projects.sh
\end{lstlisting}

\section{Company extension: processi SDLC aziendali}
\label{sec:15-company}

La company extension porta processi SDLC, standard di sicurezza e gate di governance nel contesto dell'agente.

\begin{lstlisting}[language=,caption={Struttura della company extension}]
.ai-dlc/company/
+-- README.md       ← indice
+-- PROCESS.md      ← SDLC aziendale
+-- GOVERNANCE.md   ← approvazioni, compliance
+-- STANDARDS.md    ← standard ingegneristici
+-- source/         ← PDF/DOCX originali
+-- processed/      ← versioni Markdown generate
\end{lstlisting}

Per convertire documenti aziendali in Markdown leggibile dall'agente:

\begin{lstlisting}[language=bash,caption={Preprocessing documenti aziendali}]
bash .ai-dlc/tools/preprocess-company-docs.sh
\end{lstlisting}

\begin{notebox}
Le regole della company extension hanno priorità sulle regole del framework. Se l'azienda richiede comportamenti diversi da quelli predefiniti in \texttt{.ai-dlc/modules/}, vincono le regole aziendali.
\end{notebox}

\section{Codebase legacy: analisi e documentazione}
\label{sec:15-legacy}

AI-DLC copre questo flusso con due moduli in sequenza: prima si analizza (\texttt{09\_CODEBASE\_ANALYSIS.md}), poi si documenta (\texttt{10\_DOCUMENTATION.md}).

\textbf{Step 0 --- Prepara il contesto}

\begin{lstlisting}[language=,caption={\_CONTEXT.md per analisi legacy}]
| Phase                   | 0-Discovery                               |
| Mode                    | STANDARD                                  |
| Active Task             | T-DOC-001 Analisi e doc modulo auth       |
| Active Task Model Level | 5                                         |
\end{lstlisting}

\textbf{Step 1 --- Analisi con il modulo 09}

\begin{lstlisting}[language=,caption={Prompt di analisi codebase}]
Carica .ai-dlc/modules/09_CODEBASE_ANALYSIS.md e analizza src/auth/.

Produci in docs/_analysis/:
1) MAP.md    — moduli, dipendenze, entry point
2) RISKS.md  — rischi tecnici e debito
3) HOTSPOTS.md — file più critici con motivazione
4) RUN.md    — build/test/run

Non inventare: se ti mancano informazioni, elenca ASSUMPTION + TODO.
\end{lstlisting}

\textbf{Step 2 --- Documentazione con il modulo 10}

Quando \texttt{MAP.md} è approvato, genera la documentazione navigabile usando \texttt{docs/\_analysis/} come fonte.

\textbf{Step 3 --- Verifica}

Per ogni file generato: verifica che gli endpoint API esistano nel codice, che le entità corrispondano allo schema, che i diagrammi Mermaid renderizzino e che ogni TODO e ASSUMPTION sia giustificato.

\begin{warnbox}
Non saltare l'analisi (Step 1). Documentare senza \texttt{MAP.md} produce contenuti inventati. Non usare Mode LITE o RAPID --- la documentazione è output ad alto impatto a lungo termine.
\end{warnbox}

\section*{\LabelSummary}

\begin{itemize}
  \item \textbf{Monorepo:} ogni sottoprogetto ha il proprio \texttt{\_CONTEXT.md}; l'agente usa quello più vicino al file su cui lavora.
  \item \textbf{Company extension:} \texttt{.ai-dlc/company/} porta i processi SDLC aziendali nel contesto; priorità sulle regole del framework.
  \item \textbf{Codebase legacy:} il flusso è sempre analisi (09) $\to$ documentazione (10). Non invertire l'ordine.
\end{itemize}
